\clearpage
\section{Positronium}
\subsection{Reduced Mass}
by definition
\begin{equation}
    \mu = \frac{m^2}{2m} = \frac{m}{2}
\end{equation}
\subsection{Central Force Hamiltonian}
The hamiltonian is given by the center of mass momentum element and coulomb element
\begin{equation}
    H = \frac{p^2}{2\mu} - \frac{q^2}{4\pi\epsilon_0 r}
\end{equation}
\subsection{Energy Levels and Degeneracy}
The hamiltonian is structurally the same as that of the hydrogen atom with the replacement of the mass so the spectrum is
\begin{equation}
    E_n = -\frac{\mu c^2 \alpha^2}{2n^2} = -\frac{mc^2 \alpha^2}{4n^2}
\end{equation}
which is half that of the regular hydrogen arom, and the degeneracy is the same as the hudrogen atom, \(n^2\).
\subsection{Bohr Radius}
The Bohr Radius is
\begin{equation}
    a_0 = \frac{4\pi\epsilon_0 \hbar^2}{e^2 \mu} = \frac{2\pi\epsilon_0 \hbar^2}{e^2 m} 
\end{equation}
which is double that of the hydrogen atom.
\subsection{Normalized Ground State}
The normalized wavefunction in position representation, given by spherical harmonics and Laguerre polynomials, inserting \(n=1, l=0, m=0\) is
\begin{equation}
    \psi_{100}(\vec{r}) = \sqrt{\frac{4}{a_0^3}} e^{-r/a_0} \cdot 1 \cdot \frac{1}{2}\sqrt{\frac{1}{\pi}} = \sqrt{\frac{1}{\pi a_0^3}} e^{-r/a_0} 
\end{equation}
inserting the bohr radius we found
\begin{equation}
    \left\lvert\psi_{100}(\vec{r})\right\rvert^2 = \frac{1}{\pi \ins{\frac{2\pi\epsilon_0 \hbar^2}{e^2 m} }^3} e^{-2r/\ins{\frac{2\pi\epsilon_0 \hbar^2}{e^2 m} }} 
\end{equation}
I don't see why the question asks about \(\alpha\), it doesn't depend on the mass, so it doesn't change.
\subsection{Degeneracy of Orbital Ground State}
Each of the particles in the atom has 2 possible spin states, so the total number of states is 4, which is the degeneracy.
\subsection{Gyromagnetic Ratio}
The gyromagnetic ratio is defined as 
\begin{equation}
    \gamma = \frac{q}{m}
\end{equation}
both particles' mass is the same but they have opposite charge, so \(\gamma_1 = -\gamma_2\).
\subsection{Eigenstates And Eigenvalues}
We can rewrite
\begin{equation}
    H_{SS} = \frac{A}{\hbar^2}\vec{S}\cdot\vec{S} = \frac{A}{2\hbar^2}\ins{S^2-S_1^2-S_2^2}
\end{equation}
therefore it is diagonal in the coupled basis and the eigenstates are the triplet and singlet states
\begin{gather}
    \ket{S=1, M=1} = \ket{++} \\
    \ket{S=1, M=0} = \frac{1}{\sqrt{2}}\ins{\ket{+-}+\ket{-+}} \\
    \ket{S=1, M=-1} = \ket{--} \\
    \ket{S=0, M=0} = \frac{1}{\sqrt{2}}\ins{\ket{+-}-\ket{-+}}
\end{gather}
and the eigenvalues are
\begin{gather}
    H_{SS}\ket{S=1, M=1} = \frac{A}{2}\ins{2-\frac{3}{4}-\frac{3}{4}}\ket{S=1, M=1} = \frac{A}{4}\ket{S=1, M=1}
\end{gather}
but since this is true for any \(M\),
\begin{gather}
    H_{SS}\ket{S=1, M} = \frac{A}{4}\ket{S=1, M}
\end{gather}
and
\begin{equation}
    H_{SS}\ket{S=0, M=0} = \frac{A}{2}\ins{0-\frac{3}{4}-\frac{3}{4}}\ket{S=0, M=0}  = -\frac{3A}{4}\ket{S=0, M=0} 
\end{equation}
\subsection{Difference In Hyperfine Energy}
The difference in energy between the two spin states is exactly \(A\), which itself is proportional to the multiple of the gyromagnetic ratios of both particles and the reduced mass to the 6th power (\(a_0^3\) and it contains \(\mu^2\)) 
\begin{equation}
    \Delta E \propto \gamma_1 \gamma_2 \mu^3
\end{equation}
\(\gamma_1\) is the electron's gyromagnetic ratio instead of the electrons, which is 3 orders of magnitude larger. The reduced mass is half the electron's mass. If we sum both contributions we find that the energy difference is larger by 2 orders of magnitude, which is the difference in frequencies.
\subsection{A in terms of constants}
We know that
\begin{equation}
    A = -\frac{2}{3}\frac{\hbar^2}{\epsilon_0 c^2}\gamma_1 \gamma_2 \lvert\psi_{100}(0)\rvert^2
\end{equation}
if we insert the gyromagnetic ratios and the wavefunction
\begin{align}
    A &= -\frac{2}{3}\frac{\hbar^2}{\epsilon_0 c^2}\frac{e^2}{m^2}\frac{1}{\pi a_0^3} = -\frac{2\hbar^2 e^2}{3\pi\epsilon_0 m^2 c^2}\frac{1}{\ins{\frac{2\pi\epsilon_0 \hbar^2}{e^2 m} }^3} \\
    &= -\frac{e^8 m }{12\pi^4\epsilon_0^4 c^2 \hbar^4} = \frac{1}{3}mc^2 \alpha^4
\end{align}
inserting the numbers we find
\begin{equation}
    A = \frac{1}{3}\cdot\qty{0.511}{\mega\electronvolt}\cdot \frac{1}{137^4} = \qty{4.8e-4}{\electronvolt}
\end{equation}
\subsection{Hyperfine Transition Frequency}
The hyperfine transition frequency is defined as
\begin{equation}
    f = \frac{\Delta E}{h} = \frac{A}{h} = \qty{116.1}{\giga\hertz}
\end{equation}
\subsection{Annihilation Energies}
Since the annihilation element doesn't affect the states for \(S=0\) the energy remains the same
\begin{equation}
    E_{S=0} = -\frac{3}{4}A
\end{equation}
but the energy for \(S=1\) changes to
\begin{equation}
    E_{S=1} = \frac{1+3}{4}A = A
\end{equation}
\subsection{Annihilation hyperfine frequency}
same as we did before
\begin{equation}
    f = \frac{\Delta E}{h} = \frac{7A}{4h} = \qty{203.1}{\giga\hertz}
\end{equation}